% chap07_oop.tex - Object-oriented programming basics
% ============================================================================

\chapter{Object-Oriented Programming}
\label{chap:oop}

You've been using objects throughout this book. Strings have methods like \py{.upper()}. Lists have methods like \py{.append()}. Now it's time to understand what objects are and how to create your own. Object-oriented programming (OOP) lets you bundle data and the functions that operate on that data into cohesive units.

% ----------------------------------------------------------------------------
\section{What Are Objects?}
\label{sec:what-are-objects}

An \emph{object} combines data (called \emph{attributes}) with functions that operate on that data (called \emph{methods}). You've used many objects already:

\begin{lstlisting}
text = "hello"       # text is a string object
text.upper()         # upper() is a method of string objects

numbers = [1, 2, 3]  # numbers is a list object
numbers.append(4)    # append() is a method of list objects
\end{lstlisting}

A \emph{class} is the blueprint for creating objects. \py{str} is a class; \py{"hello"} is an object (or \emph{instance}) of that class.

% ----------------------------------------------------------------------------
\section{Defining a Class}
\label{sec:defining-class}

Create a class with the \py{class} keyword:

\begin{lstlisting}
class Element:
    """Represents a chemical element."""

    def __init__(self, symbol, name, atomic_number):
        """Initialize the element with its properties."""
        self.symbol = symbol
        self.name = name
        self.atomic_number = atomic_number

    def describe(self):
        """Print a description of the element."""
        print(f"{self.name} ({self.symbol}): Z = {self.atomic_number}")
\end{lstlisting}

Let's break this down:

\begin{itemize}
    \item \py{class Element:} defines a new class called \py{Element}
    \item \py{__init__} is a special method called when you create a new object (the \emph{constructor})
    \item \py{self} refers to the object being created or used
    \item \py{self.symbol} creates an attribute called \py{symbol} on the object
    \item \py{describe} is a regular method that operates on the object
\end{itemize}

\subsection{Creating Objects}

Create objects by calling the class like a function:

\begin{lstlisting}
carbon = Element("C", "Carbon", 6)
oxygen = Element("O", "Oxygen", 8)

print(carbon.symbol)       # C
print(oxygen.atomic_number)  # 8

carbon.describe()  # Carbon (C): Z = 6
oxygen.describe()  # Oxygen (O): Z = 8
\end{lstlisting}

Each object has its own data. \py{carbon.symbol} is \py{"C"}, while \py{oxygen.symbol} is \py{"O"}.

% ----------------------------------------------------------------------------
\section{The self Parameter}
\label{sec:self}

Every method's first parameter is \py{self}, which refers to the object the method is called on:

\begin{lstlisting}
class Counter:
    def __init__(self):
        self.count = 0

    def increment(self):
        self.count += 1

    def get_count(self):
        return self.count

c = Counter()
c.increment()  # self refers to c
c.increment()
print(c.get_count())  # 2
\end{lstlisting}

When you call \py{c.increment()}, Python automatically passes \py{c} as the \py{self} parameter.

\begin{note}
The name \py{self} is a convention, not a requirement. You could use any name, but \py{self} is universal in Python---use it.
\end{note}

% ----------------------------------------------------------------------------
\section{Attributes and Methods}
\label{sec:attributes-methods}

\subsection{Instance Attributes}

Attributes defined on \py{self} belong to each individual object:

\begin{lstlisting}
class Sample:
    def __init__(self, name, concentration):
        self.name = name
        self.concentration = concentration

sample1 = Sample("A", 0.1)
sample2 = Sample("B", 0.2)

print(sample1.concentration)  # 0.1
print(sample2.concentration)  # 0.2
\end{lstlisting}

\subsection{Class Attributes}

Attributes defined on the class are shared by all objects:

\begin{lstlisting}
class Element:
    # Class attribute - shared by all instances
    elements_created = 0

    def __init__(self, symbol, name):
        self.symbol = symbol  # instance attribute
        self.name = name
        Element.elements_created += 1

e1 = Element("H", "Hydrogen")
e2 = Element("He", "Helium")
print(Element.elements_created)  # 2
\end{lstlisting}

\subsection{Methods}

Methods are functions that belong to a class:

\begin{lstlisting}
class Solution:
    def __init__(self, solute, concentration, volume):
        self.solute = solute
        self.concentration = concentration  # mol/L
        self.volume = volume  # L

    def moles(self):
        """Calculate moles of solute."""
        return self.concentration * self.volume

    def dilute(self, new_volume):
        """Return new concentration after dilution."""
        moles = self.moles()
        return moles / new_volume

    def describe(self):
        print(f"{self.concentration} M {self.solute}, {self.volume} L")


stock = Solution("NaCl", 1.0, 0.5)
stock.describe()  # 1.0 M NaCl, 0.5 L
print(f"Moles: {stock.moles()}")  # Moles: 0.5

new_conc = stock.dilute(2.0)
print(f"After dilution to 2L: {new_conc} M")
\end{lstlisting}

% ----------------------------------------------------------------------------
\section{Special Methods}
\label{sec:special-methods}

Python uses methods with double underscores (``dunder'' methods) for special behaviors:

\subsection{\_\_str\_\_ and \_\_repr\_\_}

Control how objects are displayed:

\begin{lstlisting}
class Element:
    def __init__(self, symbol, name, atomic_number):
        self.symbol = symbol
        self.name = name
        self.atomic_number = atomic_number

    def __str__(self):
        """Human-readable string (used by print)."""
        return f"{self.name} ({self.symbol})"

    def __repr__(self):
        """Developer-readable string (used in REPL)."""
        return f"Element('{self.symbol}', '{self.name}', {self.atomic_number})"

carbon = Element("C", "Carbon", 6)
print(carbon)  # Carbon (C)
\end{lstlisting}

\subsection{Comparison Methods}

Make objects comparable:

\begin{lstlisting}
class Element:
    def __init__(self, symbol, atomic_number):
        self.symbol = symbol
        self.atomic_number = atomic_number

    def __eq__(self, other):
        """Equal if same atomic number."""
        return self.atomic_number == other.atomic_number

    def __lt__(self, other):
        """Less than based on atomic number."""
        return self.atomic_number < other.atomic_number

carbon = Element("C", 6)
oxygen = Element("O", 8)

print(carbon < oxygen)   # True
print(carbon == oxygen)  # False

elements = [oxygen, carbon]
print(sorted(elements))  # sorts by atomic number
\end{lstlisting}

\subsection{Arithmetic Methods}

Define what \py{+}, \py{-}, \py{*} do for your objects:

\begin{lstlisting}
class Vector:
    def __init__(self, x, y, z):
        self.x = x
        self.y = y
        self.z = z

    def __add__(self, other):
        return Vector(self.x + other.x, self.y + other.y, self.z + other.z)

    def __mul__(self, scalar):
        return Vector(self.x * scalar, self.y * scalar, self.z * scalar)

    def __str__(self):
        return f"({self.x}, {self.y}, {self.z})"

v1 = Vector(1, 2, 3)
v2 = Vector(4, 5, 6)
print(v1 + v2)    # (5, 7, 9)
print(v1 * 2)     # (2, 4, 6)
\end{lstlisting}

% ----------------------------------------------------------------------------
\section{Encapsulation}
\label{sec:encapsulation}

\emph{Encapsulation} means hiding internal details and providing a clean interface:

\begin{lstlisting}
class Thermometer:
    def __init__(self):
        self._temperature_c = 0  # underscore suggests "private"

    def set_celsius(self, temp):
        if temp < -273.15:
            raise ValueError("Temperature below absolute zero!")
        self._temperature_c = temp

    def get_celsius(self):
        return self._temperature_c

    def get_kelvin(self):
        return self._temperature_c + 273.15

    def get_fahrenheit(self):
        return self._temperature_c * 9/5 + 32
\end{lstlisting}

The underscore prefix (\py{_temperature_c}) is a convention meaning ``don't access this directly.'' Users should use the methods, which can validate input.

\begin{note}
Python doesn't enforce privacy---you can still access \py{_temperature_c}. The underscore is a polite suggestion, not a lock.
\end{note}

% ----------------------------------------------------------------------------
\section{Inheritance}
\label{sec:inheritance}

\emph{Inheritance} lets you create new classes based on existing ones:

\begin{lstlisting}
class Compound:
    """Base class for chemical compounds."""

    def __init__(self, formula, molar_mass):
        self.formula = formula
        self.molar_mass = molar_mass

    def describe(self):
        print(f"{self.formula}: {self.molar_mass} g/mol")


class Acid(Compound):
    """An acid compound with pKa."""

    def __init__(self, formula, molar_mass, pKa):
        super().__init__(formula, molar_mass)  # call parent's __init__
        self.pKa = pKa

    def describe(self):
        super().describe()  # call parent's describe
        print(f"  pKa = {self.pKa}")

    def Ka(self):
        return 10 ** (-self.pKa)


acetic = Acid("CH3COOH", 60.05, 4.76)
acetic.describe()
# CH3COOH: 60.05 g/mol
#   pKa = 4.76

print(f"Ka = {acetic.Ka():.2e}")  # Ka = 1.74e-05
\end{lstlisting}

\py{Acid} inherits from \py{Compound}:
\begin{itemize}
    \item It gets \py{formula} and \py{molar\_mass} from \py{Compound}
    \item \py{super()} calls the parent class's methods
    \item It adds its own \py{pKa} attribute and \py{Ka()} method
    \item It can override methods (like \py{describe})
\end{itemize}

% ----------------------------------------------------------------------------
\section{Decorators and Dataclasses}
\label{sec:decorators-dataclasses}

When reading Python code, you'll encounter the \py{@} symbol above function or class definitions. This is \emph{decorator} syntax---a way to modify or enhance functions and classes.

\subsection{What Are Decorators?}

A decorator is a function that takes another function (or class) and returns a modified version of it. The \py{@decorator} syntax is shorthand:

\begin{lstlisting}
@some_decorator
def my_function():
    pass

# This is equivalent to:
def my_function():
    pass
my_function = some_decorator(my_function)
\end{lstlisting}

You don't need to write your own decorators, but you'll \emph{use} many from libraries. Common examples include \py{@property}, \py{@staticmethod}, and \py{@classmethod} in classes, plus countless decorators in frameworks like Flask or pytest.

\subsection{The @dataclass Decorator}

Python's \py{dataclasses} module provides a decorator that automatically generates common methods for classes that primarily store data:

\begin{lstlisting}
from dataclasses import dataclass

@dataclass
class Element:
    symbol: str
    name: str
    atomic_number: int
    atomic_mass: float

# The decorator automatically creates:
# - __init__ (so you can write Element("C", "Carbon", 6, 12.011))
# - __repr__ (nice string representation)
# - __eq__ (comparison by values)

carbon = Element("C", "Carbon", 6, 12.011)
print(carbon)
# Element(symbol='C', name='Carbon', atomic_number=6, atomic_mass=12.011)

oxygen = Element("O", "Oxygen", 8, 15.999)
print(carbon == oxygen)  # False
\end{lstlisting}

Compare this to the manual version, which requires much more code:

\begin{lstlisting}
class Element:
    def __init__(self, symbol, name, atomic_number, atomic_mass):
        self.symbol = symbol
        self.name = name
        self.atomic_number = atomic_number
        self.atomic_mass = atomic_mass

    def __repr__(self):
        return (f"Element(symbol='{self.symbol}', name='{self.name}', "
                f"atomic_number={self.atomic_number}, "
                f"atomic_mass={self.atomic_mass})")

    def __eq__(self, other):
        return (self.symbol == other.symbol and
                self.name == other.name and
                self.atomic_number == other.atomic_number and
                self.atomic_mass == other.atomic_mass)
\end{lstlisting}

\begin{labexample}
Dataclasses are perfect for scientific metadata:

\begin{lstlisting}
from dataclasses import dataclass
from datetime import datetime

@dataclass
class Measurement:
    sample_id: str
    value: float
    uncertainty: float
    timestamp: datetime
    instrument: str = "UV-Vis"  # default value

reading = Measurement(
    sample_id="S001",
    value=0.452,
    uncertainty=0.003,
    timestamp=datetime.now()
)
print(reading)
\end{lstlisting}
\end{labexample}

\subsection{Dataclasses vs Regular Classes}

Use a \textbf{dataclass} when:
\begin{itemize}
    \item The class primarily holds data
    \item You want automatic \py{__init__}, \py{__repr__}, \py{__eq__}
    \item Behavior is minimal or straightforward
\end{itemize}

Use a \textbf{regular class} when:
\begin{itemize}
    \item You need complex initialization logic
    \item The class has significant custom behavior
    \item You need fine control over special methods
\end{itemize}

\begin{tip}
When in doubt, start with a dataclass. You can always convert to a regular class later if you need more control. Dataclasses reduce boilerplate and make your intent clear: ``this class is primarily a data container.''
\end{tip}

% ----------------------------------------------------------------------------
\section{When to Use Classes}
\label{sec:when-classes}

Classes are useful when:

\begin{itemize}
    \item You have data that naturally belongs together (element properties, sample metadata)
    \item You need multiple instances with the same structure
    \item You have functions that always operate on the same data
    \item You want to model real-world entities (samples, instruments, experiments)
\end{itemize}

Don't use classes when:

\begin{itemize}
    \item A simple function would do
    \item You only need one instance
    \item The data doesn't have associated behavior
\end{itemize}

\begin{labexample}
A class for spectroscopy data:

\begin{lstlisting}
class Spectrum:
    def __init__(self, wavelengths, absorbances, name=""):
        self.wavelengths = wavelengths
        self.absorbances = absorbances
        self.name = name

    def peak_wavelength(self):
        """Find wavelength of maximum absorbance."""
        max_idx = self.absorbances.index(max(self.absorbances))
        return self.wavelengths[max_idx]

    def absorbance_at(self, wavelength):
        """Get absorbance at specific wavelength (nearest point)."""
        closest_idx = min(range(len(self.wavelengths)),
                         key=lambda i: abs(self.wavelengths[i] - wavelength))
        return self.absorbances[closest_idx]

    def normalize(self):
        """Return new spectrum normalized to max = 1."""
        max_abs = max(self.absorbances)
        normalized = [a / max_abs for a in self.absorbances]
        return Spectrum(self.wavelengths.copy(), normalized, self.name)


# Using the class
wl = [400, 420, 440, 460, 480, 500]
ab = [0.1, 0.3, 0.8, 0.6, 0.4, 0.2]

spectrum = Spectrum(wl, ab, "Sample A")
print(f"Peak at {spectrum.peak_wavelength()} nm")
print(f"A at 450nm: {spectrum.absorbance_at(450):.2f}")

normalized = spectrum.normalize()
print(f"Normalized peak: {max(normalized.absorbances)}")
\end{lstlisting}
\end{labexample}

% ----------------------------------------------------------------------------
\section{Chapter Summary}
\label{sec:oop-summary}

Object-oriented programming bundles data with behavior:

\begin{itemize}
    \item \textbf{Class:} Blueprint for objects (\py{class MyClass:})
    \item \textbf{Object:} Instance of a class (\py{obj = MyClass()})
    \item \textbf{Attributes:} Data stored on objects (\py{self.data})
    \item \textbf{Methods:} Functions belonging to objects
    \item \textbf{\_\_init\_\_:} Constructor, called when creating objects
    \item \textbf{self:} Reference to the current object
    \item \textbf{Inheritance:} Creating classes based on other classes
\end{itemize}

\begin{ponder}
\begin{enumerate}
    \item What's the difference between a class and an object?
    \item Why does every method need \py{self} as its first parameter?
    \item When would you use inheritance?
    \item Design a class to represent a laboratory sample with properties like ID, collection date, and concentration. What methods might it have?
\end{enumerate}
\end{ponder}
